\begin{abstract}
\noindent\textbf{Question.} Thin composite panels harden as they vibrate: their resonance
drifts upward with amplitude and, under forcing, folds and jumps. We ask whether the
\emph{spatial distribution} of fiber orientation, the freedom that automated fiber
placement provides, can reduce this hardening without changing the panel's linear
frequency or mass. If it can, we ask what mechanism provides the lever, what bounds it,
and how it compares with a uniform ply angle and with the stacking sequence.

\noindent\textbf{Hypothesis.} The linear frequency and the hardening integrate the same
rotated stiffness against two different strain fields of the fundamental: its bending
curvature and the in-plane-relaxed membrane stretch it induces. The two differ from point
to point in magnitude and principal direction. A fiber field can then take, point by point, the
direction that serves both under the frequency constraint (across the stretch, and
toward or away from the curvature as the constraint requires), while a uniform layup can
only compromise.

\noindent\textbf{Analysis.} We make this exact. A local Fourier representation shows that
each energy density depends on the fiber angle only through the anisotropic invariants
$U_2,U_3$ and the local principal directions. This predicts no lever for isotropic
material, exact invariance of a single-ply beam under any uniform angle, and a pointwise
optimality condition whose violation guarantees that steering beats every uniform angle.
We test the predictions with a spectral-submanifold reduction of laminated von K\'arm\'an
beams and plates. The hardening coefficient is exact in the in-plane relaxation and
evaluated on a mode-adaptive fundamental with an analytic adjoint. The frequency is
imposed as an exact equality and the tow-curvature limit pointwise.

\noindent\textbf{Results.}
On a clamped beam, symmetric steering with no bending--extension coupling cuts the
detuning by up to $\nBInfHeldRed\%$ with the frequency kept within $\nBeamFreqBand\%$ of the unidirectional value ($\omega_1\geq0.98\,\omega_{1,\UD}$ in the
design model; $\nBInfHeldf$ on the full model), by placing compliant fibers at the
inflection points. An optimized bending--extension coupling, the intuitive alternative, gives only $\nBBred\%$ (on its own drift objective). On a simply
supported plate, a curvature-feasible tow field reduces the frequency-normalized
hardening to $\nFTen$ of the unidirectional value at the unidirectional frequency
(held exactly in the design model, to within $\nHeldTol\%$ on refined bases and an independent finite-element model). No uniform layup can do this: the unidirectional plate is the
lowest-frequency member of all $\nSnum$ symmetric stacking sequences, and the steered
designs also undercut the best stacks at every higher frequency tested. A pointwise design computed from the
unidirectional fields alone already captures most of the gain, and steering raises the relaxed share
of the stretch energy from $\nRelaxUD\%$ to $\nRelaxOpt\%$. The lever grows monotonically
with lamina anisotropy and saturates once the minimum tow radius is below about $\SI{20}{mm}$, a fraction
of the scale of the strain fields. The full multi-mode model confirms the reduction by its SSM and by
EPMC continuation. Direct time integration shows the consequence: the jump-free
amplitude widens by $\nCacRatio\times$, as predicted by $a_c\propto J^{-1/2}$, and at a
drive where the unidirectional panel jumps with hysteresis, the steered panel is
single-valued.

\noindent\textbf{Conclusion.} Spatial fiber steering is a genuine, mechanistically
understood lever on nonlinear stiffness at fixed linear frequency and mass. It is set
by material anisotropy and by the tow radius relative to the strain-field length scale,
and for this plate it is far stronger than the through-thickness freedom of stacking.
\end{abstract}

\noindent\textbf{Keywords:} variable-stiffness composites; fiber steering; geometric nonlinearity;
spectral submanifolds; backbone curve; adjoint optimization; automated fiber placement.

